\section{生产者机制与幂等性保证}
% 负责人：范子成
% 内容：生产者机制与Kafka的幂等性
\hspace{2em}本章节主要围绕Kafka Producer的工作原理展开，重点介绍生产者发送消息的流程及其可靠性保证机制。随后，在介绍了为何引入幂等性机制后深入探讨Kafka如何通过引入Producer ID和Sequence Number来实现消息发送的幂等性，确保在网络异常或重试情况下不会产生重复消息。最后，分析幂等性机制的局限性及其适用范围。
\subsection{生产者工作原理}
% 生产者的消息发送流程、分区选择策略
    \subsubsection{生产者消息发送流程}
    
    \subsubsection{分区选择策略}

\subsection{消息发送的可靠性保证}
% acks参数、重试机制、超时处理
    \subsubsection{acks参数}

    \subsubsection{重试机制}

    \subsubsection{超时处理}

\subsection{幂等性问题的产生}
% 什么是幂等性、为什么需要幂等性
    \subsubsection{幂等性的定义}

    \subsubsection{重复发送消息的问题}
\subsection{Kafka的幂等性实现机制}
% Producer ID、Sequence Number、幂等性的实现原理
    \subsubsection{实现幂等性的相关机制}

    \subsubsection{ACK机制对比实验}

\subsection{幂等性的局限性}
% 单分区幂等性、会话级别的限制
    \subsubsection{单分区幂等性}

    \subsubsection{会话级别限制}
    
\subsection{小结}
% 本章总结
